\chapter{基于状态空间建模的脑电超分辨率方法}
\section{引言}
如第二章所述，信号级生成建模以建模观测分布$P(X)$为核心目标，通过率--失真理论指导下的变分推断框架，系统性地解决去噪去伪影、缺失插值、通道重建、统计标准化等关键问题。在多源异构EEG数据的现实场景中，不同采集设备的电极密度差异导致空间信息获取能力的显著不均衡，这种硬件层面的异构性直接影响了跨设备数据整合和知识迁移的效果。因此，本章聚焦于信号级建模中的核心问题——空间超分辨率，探讨如何通过生成式建模方法弥合低密度与高密度EEG系统之间的信息鸿沟，实现从有限空间采样到高保真度空间表征的智能重建，为后续的表征学习与多源数据融合提供统一的高质量数据基础。


\begin{figure}[htbp]
  \centering
  \includegraphics[width=\textwidth]{C3-1采集痛点.pdf}
  \caption{设备采集痛点问题}
  \caption*{\centering (a)低分辨率采集设备概览；(b)高分辨率采集设备概览。}
  \label{fig:C3-1}
\end{figure}
在现实应用中，EEG数据采集面临着显著的空间分辨率异构性挑战\cite{burns2014brain,tasci2023epilepsy,teplan2002fundamentals}。
电极密度与空间布局受到成本效益、系统复杂度、用户佩戴舒适性以及实际部署便利性等多重现实约束的制约。虽然高密度EEG系统（通常配备64--256个电极，甚至更多）能够在一定程度上提升空间采样率，如图~\ref{fig:C3-1}(b)所示，其配备的高密度电极阵列可以捕获更细粒度的空间模式信息，但同时也带来了设备成本急剧上升、复杂的电极安装与校准流程、对专业技术人员的强依赖性以及用户佩戴舒适性的显著下降等问题。这些因素使得高密度系统主要局限于研究实验室和专业医疗机构的受控环境中，难以实现大规模普及和日常应用场景的部署。
相比之下，面向消费级市场和便携式应用的低密度EEG设备（通常配备8-32个电极）在成本控制、用户友好性和部署便利性方面具有显著优势，如图~\ref{fig:C3-1}(a)所示的简化电极配置使得EEG技术能够走出实验室，进入教育、娱乐、健康监测等日常应用领域\cite{debie2019multimodal}。但其有限的电极数量和相对稀疏的空间采样不可避免地导致重要神经活动信息的丢失和空间混叠现象的加剧。根据Nyquist-Shannon采样定理，当空间采样率低于信号最高空间频率的两倍时，将产生不可恢复的混叠失真。对于头皮EEG而言，这种信息损失不仅体现在空间分辨率的直接下降，更重要的是引发了跨设备、跨场景统计特性不一致的系统性问题，进而对后续的特征提取、模式识别和决策稳定性产生连锁性的负面影响。如何在保持低密度EEG系统实用性优势的同时，通过算法途径弥补其空间信息不足的缺陷，实现“软件补偿硬件”的空间分辨率提升，已成为推动EEG技术向更广泛应用领域扩展的关键技术瓶颈。

空间超分辨率作为解决这一问题的关键技术途径，其本质目标是学习从低密度空间观测到高密度空间表征的映射函数。这一映射过程需要在保持原有神经信息完整性的同时，合理推断缺失空间位置的电位分布，实现空间细节的智能补全。围绕空间超分辨率算法，现有研究工作沿着两条主要技术路径展开深入探索。传统的数学建模方法主要基于插值理论、稀疏表示和低秩矩阵分解等经典信号处理技术，通过构建简约的数学假设和先验约束来估计缺失的空间信息\cite{candes2012exact, poon2019multidimensional}。由于EEG信号的生成机制涉及复杂的神经网络动力学、多层生物组织的体积传导效应以及电极--头皮界面的接触特性等多重非线性因素的耦合作用，传统方法基于线性假设或简单非线性模型的数学框架往往无法充分捕获这些复杂的物理和生理映射关系，导致重建质量在实际应用中显著受限\cite{yang2021asre}。
另一方面，基于深度学习的现代方法通过构建端到端的神经网络架构，利用数据驱动的表征学习能力来自动发现低分辨率与高分辨率EEG信号之间的复杂非线性映射关系，在重建质量和适应性方面取得了显著进展，涵盖了CNN\cite{saba2020eeg}、GAN\cite{corley2018deep}、Transformer架构\cite{li2023estformer}以及各种面向跨模态融合的架构\cite{shin2024super, tang2022deep}。现有深度学习方法在处理EEG这类长时序、多通道生理信号时仍面临若干关键局限。首先是计算复杂度问题，特别是Transformer类方法的$O(L^2)$二次复杂度在处理长时间窗口和高通道数场景时导致计算资源需求呈指数级增长，其次是对大规模配对训练数据的强依赖性，最后是对EEG信号固有物理--生理特性的归纳偏置不足，使得模型难以有效利用已知的神经电生理学知识的约束条件。

基于对上述技术现状和挑战的深入分析，本章提出了一种基于状态空间模型（State Space Models，SSMs）的超分辨率重建方法——MASER（\textbf{MA}mba-based \textbf{S}up\textbf{E}r-\textbf{R}esolution for EEG）。MASER的核心思想是将从低密度到高密度EEG信号的重建过程建模为一个条件生成问题，即给定低分辨率的空间观测，在先验知识注入的原则指导下，生成尽可能接近真实高密度采集系统输出的空间细节和统计结构。MASER在该领域首次系统性地引入SSMs作为核心建模工具。SSMs作为一类经典的动态系统建模方法，通过隐状态的递归更新和观测映射处理具有时间依赖性和状态演化特性的序列数据。而最近发展的选择性机制，如Mamba架构中的选择性扫描算法，进一步增强了模型对输入相关信息的自适应选择能力，通过输入依赖的参数化使得模型能够以$O(L)$线性计算复杂度有效建模长程时间依赖关系，以匹配EEG信号长时间连续记录和多通道并行处理的本质需求。整体架构采用低分辨率特征编码器和高分辨率信号解码器的设计范式，通过分阶段的特征抽象和信号重建过程，逐步提升中间表征的语义可分离性和最终输出的时空保真度。

本章将详细阐述MASER方法的技术细节，包括状态空间建模在EEG超分辨率任务中的适配机制，以及融合卷积和注意力机制的eMamba模块设计等。通过在公开数据集上的实验，验证MASER在重建质量和跨域泛化能力方面的优势。

\section{超分辨率算法框架}
% 针对消费级脑电设备电极数量有限、空间分辨率不足的现实约束，本节提出基于状态空间模型构建EEG超分辨率重建框架MASER。MASER通过学习低分辨率到高分辨率信号间的复杂映射关系，在算法层面而非硬件层面实现空间分辨率增强，提供高质量、高信噪比的数据。该方法旨在解决硬件成本与用户体验的瓶颈，为BCI系统的普及应用奠定基础。
\subsection{问题定义与优化目标}

EEG超分辨率重建问题的核心在于从有限的空间采样中恢复高分辨率的神经电信号表示，涉及神经生理学原理和信号处理理论。从数学角度，设$\mathbf{X}_l \in \mathbb{R}^{C_l \times T}$表示低分辨率EEG信号段，其中$C_l$为低分辨率通道数，$T$为时间采样点数。超分辨率重建的目标是重建高分辨率EEG信号$\mathbf{X}_h \in \mathbb{R}^{C_h \times T}$，其中$C_h > C_l$表示高分辨率信号的通道数增加。
即，给定低分辨率信号$\mathbf{X}_l = [\mathbf{x}_1, \mathbf{x}_2, \ldots, \mathbf{x}_{C_l}]^\top$和目标高分辨率信号$\mathbf{X}_h = [\mathbf{x}_1, \mathbf{x}_2, \ldots, \mathbf{x}_{C_h}]^\top$，其中$\mathbf{x}_i \in \mathbb{R}^{T}$表示第$i$个通道的时间序列数据，超分辨率过程为寻找最优映射函数$\mathcal{F}: \mathbb{R}^{C_l \times T} \rightarrow \mathbb{R}^{C_h \times T}$，使得：
\begin{equation}
\mathbf{X}_h = \mathcal{F}(\mathbf{X}_l) + \varepsilon
\label{eq:sr_mapping}
\end{equation}
式中，$\varepsilon$表示重建误差项，理想情况下应最小化其影响。根据EEG信号的信息结构定义，该映射函数需要满足多重约束条件：（1）时域连续性约束确保重建信号在时间维度上保持平滑性；（2）空间一致性约束保证重建信号符合大脑电活动的空间传播规律；（3）频域特性约束维持EEG信号的频谱分布特征；（4）生理合理性约束确保重建信号的幅值范围和统计特性符合神经生理学规律。


EEG超分辨率重建问题面临多个层面的挑战。首先是信息恢复挑战，低分辨率采样导致空间高频成分的丢失，如何从有限信息中恢复完整的空间分布是关键难题。EEG信号的空间分布遵循复杂的神经生理学规律，不同脑区之间存在功能性和解剖学联系，这些关系在低分辨率采样中可能丢失或模糊。其次是空间依赖关系建模挑战，EEG信号的空间分布不是简单的线性插值关系，而是反映了大脑皮层的电活动传播模式。这种物理约束使得相邻电极间存在复杂的非线性相关性，传统的插值方法难以准确建模这种关系。最后是计算复杂度与精度的平衡挑战，深度学习方法如Transformer在处理长序列数据时面临二次复杂度增长的问题，而EEG信号通常包含数千个时间点，这使得计算成本急剧上升。

为了有效解决这些挑战，MASER采用基于状态空间模型的深度学习框架。该方法的核心思想是将EEG超分辨率问题转化为一个条件生成建模问题，通过学习低分辨率到高分辨率信号之间的复杂映射关系，实现高质量的空间分辨率增强。MASER将输入的低分辨率信号视为观测序列，通过状态空间模型的动态演化机制捕获信号的时空依赖关系，进而重建高分辨率信号。

\subsection{状态空间模型基础原理}
SSMs源自控制理论与信号处理，旨在以一组内部潜在变量刻画受输入驱动并随时间演化的动态系统。
这一建模范式对于EEG信号尤为契合，因为其记录到的头皮电位是多源神经元群同步活动经体积传导后的混合投影，潜在的神经动力学状态不可直接观测但能够通过观测方程间接推断 \cite{teplan2002fundamentals, ferree2001spatial, friston2008multiple}。在此框架中，动力学与观测两条链路相互耦合，可统一容纳系统输入、内部状态演化、以及多通道观测的生成机制。


在连续时间域，经典线性时不变SSMs可由如下耦合微分方程给出：
\begin{equation}
\left\{
\begin{aligned}
\dot{\mathbf{x}}(t) &= \mathbf{A}\mathbf{x}(t) + \mathbf{B}\mathbf{u}(t), \\
\mathbf{y}(t) &= \mathbf{C}\mathbf{x}(t) + \mathbf{D}\mathbf{u}(t),
\end{aligned}
\right.
\label{eq:continuous_ssm}
\end{equation}
式中，$\mathbf{x}(t)\in\mathbb{R}^{n}$表示时刻$t$的潜在状态向量，$\mathbf{u}(t)\in\mathbb{R}^{m}$表示外部输入（或刺激）向量，$\mathbf{y}(t)\in\mathbb{R}^{p}$表示观测向量；$\mathbf{A}\in\mathbb{R}^{n\times n}$表示状态转移矩阵，描述状态的内在动力学；$\mathbf{B}\in\mathbb{R}^{n\times m}$表示输入矩阵，描述外部输入对状态的作用；$\mathbf{C}\in\mathbb{R}^{p\times n}$表示观测矩阵，将潜在状态映射到可观测空间；$\mathbf{D}\in\mathbb{R}^{p\times m}$表示直接传输矩阵，刻画输入对观测的瞬时影响。
在EEG语境中，$\mathbf{u}(t)$可对应感觉或任务刺激、药物干预或外源噪声干扰，$\mathbf{C}$体现体积传导与导联布局对潜在皮层源活动的线性投影\cite{ferree2001spatial}。当$\mathbf{A}$的特征值位于复平面左半侧时，系统在连续时间上是渐近稳定的，这与许多生理过程的耗散特性相符。

深度学习框架通常处理离散时间序列，需对式\eqref{eq:continuous_ssm}进行离散化。采用零阶保持（Zero-Order Hold, ZOH）假设在采样周期内输入保持常值\cite{galias2008analysis, gupta2022diagonal}，可得：
\begin{equation}
\left\{
\begin{aligned}
\mathbf{x}[k+1] &= \mathbf{A}_d\mathbf{x}[k] + \mathbf{B}_d\mathbf{u}[k] + \mathbf{w}[k], \\
\mathbf{y}[k] &= \mathbf{C}_d\mathbf{x}[k] + \mathbf{D}_d\mathbf{u}[k] + \mathbf{v}[k],
\end{aligned}
\right.
\label{eq:discrete_ssm}
\end{equation}
式中，$k$表示离散时间步，$\Delta t$为采样周期，$\mathbf{A}_d=e^{\mathbf{A}\Delta t}$，$\mathbf{B}_d=\int{0}^{\Delta t}e^{\mathbf{A}\tau}\mathbf{B},d\tau$，$\mathbf{C}_d=\mathbf{C}$，$\mathbf{D}_d=\mathbf{D}$，$\mathbf{w}[k]$与$\mathbf{v}[k]$分别为离散过程与观测噪声，协方差记为$\mathbf{Q}_d$与$\mathbf{R}_d$。$\mathbf{Q}_d$与连续协方差$\mathbf{Q}_c$满足矩阵Lyapunov型关系，而$\mathbf{R}_d=\mathbf{R}_c$在理想瞬时采样时近似成立。式\eqref{eq:discrete_ssm}中的矩阵指数与积分项可通过分块矩阵指数或Krylov方法高效求解\cite{boley1994krylov}，在高维情形可采用对角或分块对角参数化以提升数值稳定性与训练可控性\cite{gupta2022diagonal, gu2022parameterization}。
模型辨识与推断的可行性依赖于可控性与可观性：对$(\mathbf{A},\mathbf{B})$，若可控性矩阵$\mathcal{C}=[\mathbf{B}, \mathbf{A}\mathbf{B}, \dots, \mathbf{A}^{n-1}\mathbf{B}]$满秩，则状态对输入是可达的；对$(\mathbf{A},\mathbf{C})$，若可观性矩阵$\mathcal{O}=[\mathbf{C}^{\top}, (\mathbf{C}\mathbf{A})^{\top}, \dots, (\mathbf{C}\mathbf{A}^{n-1})^{\top}]^{\top}$满秩，则潜在状态可由观测重构。

对EEG信号而言，导联数量$p$通常远小于潜在源维度$n$，因此需借助稀疏先验或低秩约束提升可观性等效度\cite{friston2008multiple}。在稳定性方面，离散谱半径$\rho(\mathbf{A}_d)<1$对应渐近稳定，这对长序列建模中的梯度传播与记忆保持同样重要\cite{gu2021efficiently}。与注意力机制在序列长度$N$上通常呈$O(N^{2})$复杂度不同，基于SSMs的前向传播可实现$O(N\cdot D)$的线性复杂度（$D$为状态维度），在长序列EEG高采样率、长时段记录的学习中具有显著效率优势\cite{gu2021efficiently, gupta2022diagonal}。此外，SSMs的卷积化实现可通过频域或递推内核加速，并借助参数化稳定化技巧（如对角稳定参数化）保证训练收敛\cite{gu2022parameterization}。
近期的选择性状态空间模型（Selective SSMs）在深度序列学习中展现出优异的长程建模能力与计算效率。例如，Mamba引入输入依赖的门控与动态滤波系数，使状态更新与读出在时刻$k$对$\mathbf{u}[k]$呈选择性调制\cite{gu2023mamba}。抽象地，可将其理解为对式\eqref{eq:discrete_ssm}中$\mathbf{A}_d,\mathbf{B}_d,\mathbf{C}_d,\mathbf{D}_d$的时变、上下文依赖参数化：
\begin{equation}
\left\{
\begin{aligned}
\mathbf{x}[k+1] &= \underbrace{\mathbf{A}_d(\mathbf{u}[k])}_{\text{选择性动力学}}\mathbf{x}[k] + \underbrace{\mathbf{B}_d(\mathbf{u}[k])}_{\text{选择性驱动}}\mathbf{u}[k], \\
\mathbf{y}[k] &= \underbrace{\mathbf{C}_d(\mathbf{u}[k])}_{\text{选择性读出}}\mathbf{x}[k] + \underbrace{\mathbf{D}_d(\mathbf{u}[k])}_{\text{选择性直达}}\mathbf{u}[k],
\end{aligned}
\right.
\label{eq:selective_ssm}
\end{equation}
式中，$\mathbf{A}_d(\cdot)$、$\mathbf{B}_d(\cdot)$、$\mathbf{C}_d(\cdot)$与$\mathbf{D}_d(\cdot)$表示由当前输入驱动的时变映射，实作上通过高效的参数化核与门控机制实现线性时间复杂度\cite{gu2023mamba, ali2024hidden, xu2024survey}。相较于固定系数的SSMs，选择性机制能够在含噪、非平稳的EEG序列中自适应地传播或遗忘信息，突出任务相关的时空模式并抑制伪迹干扰，这一思想已在计算机视觉与自然语言的数据分析中被验证\cite{liu2024vmamba, zhu2024vision, li2024mamba}。

\subsection{MASER网络架构设计}
基于前述状态空间模型理论，MASER框架构建了端到端的EEG超分辨率重建网络，整体架构如图~\ref{MASER_architecture}所示。该架构采用双阶段设计策略，通过低分辨率特征编码器与高分辨率信号解码器的协同工作，实现从有限空间采样到完整空间分布的映射重建。MASER的设计理念深度融合了神经生理学原理与深度学习技术，通过状态空间模型的动态建模能力，有效捕获EEG信号中复杂的时空依赖关系，并在计算效率与重建精度之间取得了良好平衡。

\begin{figure*}[htbp]
	\centering
		\includegraphics[width=\textwidth]{C3-2MASER框架-CN.pdf}
    \caption{EEG空间超分辨率网络MASER}
    \caption*{\centering (a)整体网络架构；(b)eMamba模块；(c)eSS2D模块。}
\label{MASER_architecture}
\end{figure*}

如算法\ref{alg:eeg_super_resolution}所述，MASER将EEG超分辨率问题转化为序列到序列的条件生成任务，通过学习低分辨率到高分辨率信号间的复杂非线性映射关系，在算法层面实现空间分辨率的增强。该方法摒弃了传统基于插值或简单卷积的信号重建策略，转而采用基于状态空间模型的深度学习框架，能够更好地建模EEG信号的动态特性和空间相关性。整个网络架构遵循编码器--解码器的设计范式，但在特征表示和序列建模方面引入了专门针对EEG信号特性优化的组件。

\begin{algorithm}[htbp]
\caption{\textbf{脑电超分辨率算法MASER流程}}
\label{alg:eeg_super_resolution}
\KwIn{低分辨率脑电信号段 $\mathbf{X}_l$: {\texttt{(C\textsubscript{l}, T)}}}
\KwOut{高分辨率脑电信号 $\mathbf{X}_h$: {\texttt{(C\textsubscript{h}, T)}}}
\SetKwComment{Comment}{$\triangleright$\ }{}

\Comment{将$\mathbf{X}$切分为长度为$T_w$的窗口并嵌入到$L$维特征中}
\Begin{
    $\mathbf{X}_{\text{patch}}$: {\texttt{(N $\times$ C\textsubscript{l}, T\textsubscript{w})}} $\leftarrow \text{Slice}(\mathbf{X}, T_w)$;
    
    $\mathbf{X}_{\text{embedded}}$: {\texttt{(N $\times$ C\textsubscript{l}, L)}} $\leftarrow \text{Embed}(\mathbf{X}_{\text{patch}}, L)$;
}

\Comment{通过堆叠的eMamba块提取低分辨率脑电特征}
$\mathbf{F}_{\text{LR}}$: {\texttt{(N $\times$ C\textsubscript{l}, L)}} $\leftarrow \text{LRExtractor}(\mathbf{X}_{\text{patch}})$;

\Comment{生成初始高分辨率占位符标记，$\Delta C = C_h - C_l$}
\Begin{
    $\mathbf{X}_{\text{mean}}$: {\texttt{(1, L)}} $\leftarrow \text{Mean}(\mathbf{X}_{\text{patch}})$;
    
    $\mathbf{X}_{\text{predict}}$: {\texttt{(N $\times$ $\Delta$C, L)}} $\leftarrow \text{Replicate}(\mathbf{X}_{\text{mean}})$;
    
    $\mathbf{F}_{\text{mean}}$: {\texttt{(1, L)}} $\leftarrow \text{Mean}(\mathbf{F}_{\text{LR}})$;
    
    $\mathbf{F}_{\text{predict}}$: {\texttt{(N $\times$ $\Delta$C, L)}} $\leftarrow \text{Replicate}(\mathbf{F}_{\text{mean}})$;
    
    $\mathbf{F}_{\text{predict}}$: {\texttt{(N $\times$ $\Delta$C, L)}} $\leftarrow \mathbf{X}_{\text{predict}} + \mathbf{F}_{\text{predict}}$;

    $\mathbf{F}_{\text{HR}}$:  {\texttt{(N $\times$ C\textsubscript{h}, L)}} $\leftarrow \text{Concatenate}(\mathbf{F}_{\text{LR}}, \mathbf{F}_{\text{predict}})$;
}

\Comment{通过堆叠的eMamba块预测高分辨率脑电信号}
\Begin{
    $\mathbf{F}_{\text{HR}}$: {\texttt{(N $\times$ C\textsubscript{h}, L)}} $\leftarrow \text{HRPredictor}(\mathbf{F}_{\text{HR}})$;
    
    $\mathbf{X}_h$: {\texttt{(C\textsubscript{h}, T)}} $\leftarrow \text{Rearrange}(\mathbf{F}_{\text{HR}})$;
}

\Return $\mathbf{X}_h$\;
\end{algorithm}


给定低分辨率输入$\mathbf{X}_l \in \mathbb{R}^{C_l \times T}$，网络首先通过分块嵌入机制将连续时间序列转换为离散的特征表示。具体而言，输入信号被划分为长度为$T_w$的时间窗口，形成分块矩阵$\mathbf{X}_{\text{patch}} \in \mathbb{R}^{(N \times C_l) \times T_w}$，其中$N = \lceil T/T_w \rceil$表示时间窗口数量。这种分块策略既保持了EEG信号的时间局部性，又为并行处理提供了便利。随后，每个时间窗口通过学习的线性变换映射到$L$维特征空间，得到嵌入表示$\mathbf{X}_{\text{embedded}} \in \mathbb{R}^{(N \times C_l) \times L}$：
\begin{equation}
\mathbf{X}_{\text{embedded}}[i] = \mathbf{W}_{\text{embed}} \mathbf{X}_{\text{patch}}[i] + \mathbf{b}_{\text{embed}},
\label{eq:patch_embedding}
\end{equation}
式中，$\mathbf{W}_{\text{embed}} \in \mathbb{R}^{L \times T_w}$表示可学习的嵌入权重矩阵，用于将时间窗口映射到特征空间，$\mathbf{b}_{\text{embed}} \in \mathbb{R}^{L}$表示偏置向量，用于增加模型的表达能力，$i \in \{1, 2, \ldots, N \times C_l\}$表示分块索引，覆盖所有通道的所有时间窗口，$L$表示嵌入特征的维度。分块嵌入策略不仅保持了时间序列的局部结构特性，还为后续的状态空间建模提供了合适的输入维度，使得模型能够在保持计算效率的同时充分利用时间域信息。

低分辨率特征编码器作为网络的前端模块，承担着从输入的低分辨率EEG信号中提取高质量时空特征表示的重要任务。该模块由多个堆叠的eMamba块组成，eMamba块整合了层归一化（Layer Normalization，LN）、EEG选择性扫描二维（EEG Selective Scan 2D，eSS2D）块以及通道注意力块（Channel Attention Block，CAB），形成了一个功能完整且高效的特征提取单元。eMamba块的处理流程体现了残差连接和注意力机制的有机结合，其数学表达式可以描述为：
\begin{equation}
\left\{
\begin{aligned}
\mathbf{Z}^{(i)} &= \mathrm{LayerNorm}\!\left(\mathbf{F}^{(i-1)}\right), \\
\mathbf{H}^{(i)} &= \mathrm{eSS2D}\!\left(\mathbf{Z}^{(i)}\right) + \mathbf{Z}^{(i)}, \\
\mathbf{F}^{(i)} &= \mathrm{CAB}\!\left(\mathbf{H}^{(i)}\right) + \mathbf{H}^{(i)},
\end{aligned}
\right.
\label{eq:emamba_block}
\end{equation}
式中，$\mathbf{F}^{(i)} \in \mathbb{R}^{(N \times C_l) \times L}$表示第$i$个eMamba块的输出特征张量，包含了经过多层处理后的丰富特征表示，$\mathbf{F}^{(i-1)} \in \mathbb{R}^{(N \times C_l) \times L}$表示前一层eMamba块的输出特征，作为当前层的输入，$\mathbf{Z}^{(i)} \in \mathbb{R}^{(N \times C_l) \times L}$表示经过层归一化处理后的中间特征，用于稳定训练过程并加速收敛，$\mathbf{H}^{(i)} \in \mathbb{R}^{(N \times C_l) \times L}$表示经过eSS2D块处理并添加残差连接后的特征表示，融合了状态空间建模的动态信息，$\mathrm{LayerNorm}(\cdot)$表示层归一化操作，用于标准化输入特征的分布，$\mathrm{eSS2D}(\cdot)$表示EEG选择性扫描二维块，实现基于状态空间模型的时空特征建模，$\mathrm{CAB}(\cdot)$表示通道注意力块，通过学习通道间的依赖关系增强重要特征。

特征扩展机制是MASER实现空间分辨率增强的关键技术组件，其设计目标是在保持原有通道信息完整性的基础上，合理地生成新增通道的初始特征表示。这一机制的核心思想是利用已有的低分辨率信息来指导高分辨率特征的初始化，避免随机初始化可能带来的训练不稳定问题。设$\Delta C = C_h - C_l$表示需要增加的通道数量，即超分辨率过程中新增的通道数目。特征扩展过程首先计算输入分块和低分辨率特征的统计信息，然后基于这些统计信息生成合理的占位符标记。
\begin{equation}
\left\{
\begin{aligned}
\mathbf{X}_{\text{mean}} &= \frac{1}{N \times C_l} \sum_{i=1}^{N \times C_l} \mathbf{X}_{\text{patch}}[i], \\
\mathbf{F}_{\text{mean}} &= \frac{1}{N \times C_l} \sum_{i=1}^{N \times C_l} \mathbf{F}_{\text{LR}}[i],
\end{aligned}
\right.
\label{eq:mean_computation}
\end{equation}
式中，$\mathbf{X}_{\text{mean}} \in \mathbb{R}^{1 \times L}$表示所有输入分块在特征维度上的平均值，反映了输入信号的整体特征分布，$\mathbf{F}_{\text{mean}} \in \mathbb{R}^{1 \times L}$表示所有低分辨率特征在特征维度上的平均值，包含了经过特征提取后的全局信息，$\mathbf{X}_{\text{patch}}[i] \in \mathbb{R}^{L}$表示第$i$个输入分块的特征向量，$\mathbf{F}_{\text{LR}}[i] \in \mathbb{R}^{L}$表示第$i$个低分辨率特征向量，$N \times C_l$表示总的分块数量，覆盖所有通道和所有时间窗口。

接下来，将这些均值表示复制$N \times \Delta C$次，为新增通道生成初始特征表示：
\begin{equation}
\left\{
\begin{aligned}
\mathbf{X}_{\text{predict}} &= \mathrm{Replicate}\!\left(\mathbf{X}_{\text{mean}}, N \times \Delta C\right), \\
\mathbf{F}_{\text{predict}}^{(\text{init})} &= \mathrm{Replicate}\!\left(\mathbf{F}_{\text{mean}}, N \times \Delta C\right), \\
\mathbf{F}_{\text{predict}} &= \mathbf{X}_{\text{predict}} + \mathbf{F}_{\text{predict}}^{(\text{init})},
\end{aligned}
\right.
\label{eq:feature_replication}
\end{equation}
式中，$\mathbf{X}_{\text{predict}} \in \mathbb{R}^{(N \times \Delta C) \times L}$表示基于输入均值生成的预测通道初始特征，$\mathbf{F}_{\text{predict}}^{(\text{init})} \in \mathbb{R}^{(N \times \Delta C) \times L}$表示基于特征均值生成的预测通道初始特征，$\mathbf{F}_{\text{predict}} \in \mathbb{R}^{(N \times \Delta C) \times L}$表示通过残差连接结合后的最终预测通道特征，$\mathrm{Replicate}(\cdot, n)$表示将输入向量复制$n$次的操作。最后，通过特征拼接操作将原有的低分辨率特征与新生成的预测特征结合：
\begin{equation}
\mathbf{F}_{\text{HR}} = \mathrm{Concatenate}\!\left(\mathbf{F}_{\text{LR}}, \mathbf{F}_{\text{predict}}\right),
\label{eq:feature_concatenation}
\end{equation}
式中，$\mathbf{F}_{\text{LR}} \in \mathbb{R}^{(N \times C_l) \times L}$表示低分辨率特征编码器的输出特征，保留了原始通道的所有信息，$\mathbf{F}_{\text{HR}} \in \mathbb{R}^{(N \times C_h) \times L}$表示扩展后的高分辨率特征，包含了原有通道和新增通道的完整特征表示，$\mathrm{Concatenate}(\cdot, \cdot)$表示沿通道维度的拼接操作。该设计既保持了原有通道的特征信息完整性，又为新增通道提供了合理的初始化表示，避免了随机初始化可能导致的训练不稳定问题。

高分辨率信号解码器作为网络的后端核心模块，承担着将扩展后的特征表示转换为最终高分辨率EEG信号的任务。该模块的设计充分考虑了特征精细化处理和信号重建的双重需求，通过多层次的特征变换和非线性映射，确保输出信号的质量和生理合理性。预测器同样采用多个eMamba块的堆叠结构，但在网络的末端添加了潜在层和多层感知机（Multi-Layer Perceptron，MLP）层，以实现从高维特征空间到原始信号空间的有效映射。整个预测过程可以分解为特征精细化、潜在表示学习和信号重建三个阶段：
\begin{equation}
\left\{
\begin{aligned}
\mathbf{F}_{\text{refined}} &= \mathrm{HRPredictor}_{\text{eMamba}}\!\left(\mathbf{F}_{\text{HR}}\right), \\
\mathbf{F}_{\text{output}} &= \mathrm{MLP}\!\left(\mathrm{LatentLayer}\!\left(\mathbf{F}_{\text{refined}}\right)\right), \\
\mathbf{X}_h &= \mathrm{Rearrange}\!\left(\mathbf{F}_{\text{output}}\right),
\end{aligned}
\right.
\label{eq:hr_prediction}
\end{equation}
式中，$\mathbf{F}_{\text{refined}} \in \mathbb{R}^{(N \times C_h) \times L}$表示经过高分辨率预测器中eMamba块处理的精细化特征，包含了经过多层非线性变换后的丰富特征表示，$\mathrm{HRPredictor}_{\text{eMamba}}(\cdot)$表示高分辨率预测器中的eMamba块堆叠结构，用于进一步提取和精细化特征，$\mathrm{LatentLayer}(\cdot)$表示潜在层操作，用于学习特征的低维潜在表示并增强模型的泛化能力，$\mathrm{MLP}(\cdot)$表示多层感知机网络，实现从潜在特征空间到输出特征空间的非线性映射，$\mathbf{F}_{\text{output}} \in \mathbb{R}^{(N \times C_h) \times T_w}$表示预测器的输出特征，具有与原始时间窗口相同的时间维度，$\mathbf{X}_h \in \mathbb{R}^{C_h \times T}$表示最终重建的高分辨率EEG信号，通过重新排列操作恢复原始的通道--时间格式，$\mathrm{Rearrange}(\cdot)$表示张量重排操作，将分块特征重新组织为连续的时间序列格式。

MASER网络架构的整体训练过程遵循端到端的优化策略，通过最小化综合损失函数来学习从低分辨率到高分辨率信号的最优映射参数。该框架的训练目标不仅包括信号重建的准确性，还考虑了生理合理性、频域特性保持以及空间一致性等多个约束条件，以确保重建信号在多个维度上都能达到高质量标准。

\subsection{eSS2D特征提取模块设计}
eSS2D模块构成了MASER架构的关键技术组件，基于视觉领域的二维选择性扫描机制发展而来，并经过深度改进以适配EEG信号的多通道时序特性和生理约束。如图~\ref{fig:C3-3}所示，该模块将多通道EEG信号的二维时空表示转换为一维序列表示，从而充分利用状态空间模型在序列建模方面的强大能力，同时通过多方向扫描策略保持原有的空间关系信息和时间依赖结构。eSS2D模块的处理流程包含分块嵌入与预处理、通道--时间增强（Channel-Temporal Enhancement，CTE）、四方向选择性扫描建模和多方向特征融合四个关键阶段。模块设计融合了EEG信号的物理特性，包括电极间的空间分布、信号的时间动态、以及不同频率成分的相互作用，确保特征提取过程能够最大程度地保留和利用这些重要信息。

\begin{figure}[htb]
\centering
\includegraphics[width=\textwidth]{C3-3eSS2D-CN.pdf}
\caption{eSS2D模块设计细节}
\label{fig:C3-3}
\end{figure}

在分块嵌入与预处理阶段，输入的低分辨率EEG信号$\mathbf{X}_l \in \mathbb{R}^{C_l \times T}$首先被分割为长度为$T_w$的时间窗口，形成分块表示$\mathbf{X}_{\text{patch}} \in \mathbb{R}^{(N \times C_l) \times T_w}$。这种分块化处理策略旨在平衡信息完整性与计算效率之间的关系。具体的分块操作可以表示为：
\begin{equation}
\mathbf{X}_{\text{patch}}[i \cdot C_l + j, :] = \mathbf{X}_l[j, i \cdot T_w : (i+1) \cdot T_w],
\label{eq:patch_division}
\end{equation}
式中，$N = \lceil T/T_w \rceil$表示总的时间窗口数量，通过向上取整确保覆盖所有时间点，$C_l$表示低分辨率EEG信号的通道数量，$T$表示信号的总时间长度，$T_w$表示每个时间窗口的长度，$i \in \{0, 1, \ldots, N-1\}$表示时间窗口索引，$j \in \{0, 1, \ldots, C_l-1\}$表示通道索引，$\mathbf{X}_{\text{patch}}[i \cdot C_l + j, :] \in \mathbb{R}^{T_w}$表示第$j$个通道第$i$个时间窗口的信号片段。分块策略能够并行处理不同的时间段，同时保持局部时间结构的完整性。窗口长度$T_w$的选择需要综合考虑多个因素：过短的窗口可能无法捕获完整的神经事件（如事件相关电位的完整波形），而过长的窗口会显著增加计算复杂度并可能引入不相关的背景噪声。


CTE阶段的设计目标是识别并突出对超分辨率重建任务最关键的时间特征成分。该阶段引入了自适应权重学习机制，能够根据输入信号的具体内容动态调整不同时间点的重要性权重，从而在后续处理中更好地保留关键信息。CTE模块通过学习通道--时间权重矩阵来实现这一目标：
\begin{equation}
\mathbf{A}_{\text{ct}} = \sigma\left(\mathrm{Conv1D}\left(\mathrm{Avg}\left(\mathbf{X}_{\text{patch}}\right)\right)\right),
\label{eq:cte_weights}
\end{equation}
式中，$\mathrm{Avg}(\cdot)$表示在$N \times C_l$维度上计算平均值的操作，将输入张量从$(N \times C_l) \times T_w$维度压缩到$1 \times T_w$维度，$\mathrm{Conv1D}(\cdot)$表示核大小为1的一维卷积操作，用于学习时间维度上的特征权重而保持输入输出维度一致，$\sigma(\cdot)$表示Sigmoid激活函数，将权重值归一化到$(0,1)$区间以确保数值稳定性，$\mathbf{A}_{\text{ct}} \in \mathbb{R}^{1 \times T_w}$表示学习得到的通道--时间权重矩阵，其每个元素对应一个时间点的重要性权重。这种设计能够自适应地学习不同时间点的重要性分布，有效反映EEG信号在不同时间段的信息密度差异和神经活动的动态变化特性。权重矩阵计算完成后，通过逐元素乘法将其应用到原始分块数据上，实现时间维度的选择性增强：
\begin{equation}
\mathbf{X}_{\text{enhanced}} = \mathbf{X}_{\text{patch}} \odot \mathbf{A}_{\text{ct}},
\label{eq:enhancement_application}
\end{equation}
式中，$\odot$表示逐元素乘法，通过广播机制作用于$N \times C_l$维度，$\mathbf{X}_{\text{enhanced}} \in \mathbb{R}^{(N \times C_l) \times T_w}$表示经过时间增强处理后的EEG分块数据，保持与输入相同的维度结构。增强后的信号在保持原有时空维度结构的同时，通过权重调制突出了时间维度上的关键信息成分，这种选择性增强机制有助于提高后续状态空间建模的效果和重建质量。

选择性扫描机制旨在通过多角度的序列化策略充分捕获二维时空数据中的结构信息和依赖关系。该机制设计了四种不同的扫描方向，每种方向都对应特定的序列化路径和建模目标：水平从左到右扫描（Horizontal Top-to-Bottom，HTB）、水平从右到左扫描（Horizontal Bottom-to-Top，HBT）、垂直从上到下扫描（Vertical Top-to-Bottom，VTB）、垂直从下到上扫描（Vertical Bottom-to-Top，VBT）。
水平扫描模式（HTB和HBT）主要用于捕获同一时间点不同通道间的空间关系和功能连接模式。这种扫描模式建模大脑皮层的横向连接、功能网络的同步活动以及不同脑区间的瞬时相互作用。通过水平方向的序列化，模型能够学习到电极空间分布中的邻近关系和远程连接模式，这对于理解大脑的网络化组织结构具有重要意义。垂直扫描模式（VTB和VBT）则建模同一通道不同时间点间的时序关系和动态演化过程，能够有效捕获神经活动的时间动态特性、振荡模式以及事件相关的时间锁定响应。对于每个扫描方向$d \in \{\text{HTB}, \text{HBT}, \text{VTB}, \text{VBT}\}$，增强后的EEG数据$\mathbf{X}_{\text{enhanced}}$按照特定的空间路径被序列化为一维表示$\mathbf{X}_d$，然后通过选择性状态空间（S6）\cite{gu2021efficiently}块进行特征交互和状态更新：
\begin{equation}
\mathbf{Y}_d = \mathbf{A}_d \mathbf{H}_d + \mathbf{B}_d \mathbf{X}_d ,
\label{eq:s6_processing}
\end{equation}
式中，$d$表示具体的扫描方向，包括HTB、HBT、VTB、VBT四种模式，$\mathbf{H}_d \in \mathbb{R}^{n \times L}$表示方向特定的初始隐状态向量，包含了状态空间模型的内部记忆信息，$\mathbf{X}_d \in \mathbb{R}^{(N \times C_l \times T_w) \times L}$表示按照方向$d$序列化后的输入特征序列，$\mathbf{A}_d \in \mathbb{R}^{n \times n}$表示方向特定的状态转移矩阵，描述系统内部状态的动态演化规律，$\mathbf{B}_d \in \mathbb{R}^{n \times L}$表示方向特定的输入矩阵，控制外部输入对内部状态的影响程度，$\mathbf{Y}_d \in \mathbb{R}^{(N \times C_l \times T_w) \times L}$表示经过状态空间建模后的输出特征序列，$n$表示隐状态的维度，$L$表示特征嵌入的维度。S6块通过递归状态更新机制建模长程时空依赖关系，同时保持线性计算复杂度的优势。不同方向的状态转移矩阵通过端到端训练自动学习，使得模型能够自适应地捕获EEG信号中复杂的方向性时空动态模式。

多方向特征融合阶段负责将四个不同扫描方向的输出特征有效整合为统一的特征表示。简单的特征拼接或算术平均操作可能会丢失重要的方向性信息和空间结构，eSS2D采用了可学习加权融合策略：
\begin{equation}
\mathbf{F} = \sum_{d \in \{\text{HTB}, \text{HBT}, \text{VTB}, \text{VBT}\}} \alpha_d \cdot \mathbf{Y}_d ,
\label{eq:weighted_fusion}
\end{equation}
式中，$\alpha_d \in \mathbb{R}^+$表示方向$d$对应的可学习融合权重，反映该方向特征对最终结果的贡献程度，$\mathbf{F} \in \mathbb{R}^{(N \times C_l \times T_w) \times L}$表示融合后的最终特征表示，包含了来自所有扫描方向的综合信息，权重参数满足归一化约束$\sum_{d} \alpha_d = 1$以确保特征融合的稳定性和可解释性。这些融合权重通过端到端的反向传播算法自动学习和优化，能够根据不同的信号特征、任务需求和数据分布动态调整各个方向的贡献度。在实际训练过程中，模型会自动学习到不同扫描方向对于特定类型EEG信号和重建任务的重要性分布，形成自适应的特征融合机制。

与二维卷积操作相比，eSS2D能够建模更长距离的时空依赖关系，不受卷积核大小的严格限制，能够处理EEG信号中常见的长程神经连接和跨时间尺度的动态过程。与Transformer架构的全局注意力机制相比，eSS2D具有线性而非二次的计算复杂度，在处理长序列EEG数据时展现出更好的计算效率和内存使用效率，同时避免了注意力机制在长序列上可能出现的梯度消失和训练不稳定问题。与循环神经网络（RNN/LSTM）相比，eSS2D通过并行化的多方向扫描策略显著提高了计算效率，同时通过状态空间模型的设计有效避免了循环网络中常见的梯度消失和梯度爆炸问题。此外，eSS2D的多方向扫描机制还提供了更好的可解释性，能够分析不同方向特征对最终重建结果的具体贡献，这对于理解模型的工作机制和优化网络结构具有重要价值。

\subsection{多组件复合损失函数}
EEG信号作为反映大脑神经活动的电生理信号，其质量评估不能仅仅依赖于简单的数值误差度量，还必须考虑信号的生理合理性、时空连续性以及频域特性的保持。单一损失函数往往难以同时满足这些多维度的质量要求，容易导致重建信号在某些方面表现良好但在其他关键特性上存在缺陷。因此，MASER采用了多个互补的损失项的加权组合，协同作用以实现全面的信号质量提升。


该损失函数的设计充分考虑了EEG信号作为生物电信号的多维特性和复杂约束条件，由均方误差损失项、时域平滑约束项和模型正则化项三个核心组件构成。这三个组件分别从重建精度、时域连续性和模型泛化能力三个维度对模型进行约束和指导。损失函数的数学表达式可以表示为：
\begin{equation}
\begin{aligned}
\mathcal{L}_{\text{SR}} &= \mathcal{L}_{\text{MSE}} + \alpha \mathcal{L}_{\text{smooth}} + \beta \mathcal{L}_{\text{reg}} \\
&= \underbrace{\frac{1}{C_{\text{up}} T} \sum_{i=1}^{C_{\text{up}}} \sum_{t=1}^{T} \left(\mathbf{Y}_{i,t} - \hat{\mathbf{Y}}_{i,t}\right)^2}_{\text{均方误差}} 
 + \alpha \underbrace{\sum_{i=1}^{C_{\text{up}}} \sum_{t=1}^{T-1} \left(\hat{\mathbf{Y}}_{i,t+1} - \hat{\mathbf{Y}}_{i,t}\right)^2}_{\text{时域平滑约束}}  + \beta \underbrace{\sum_{j=1}^{P} \|\mathbf{W}_j\|_2^2}_{\text{参数正则化}},
\end{aligned}
\end{equation}
式中，$\mathbf{Y}_{i,t} \in \mathbb{R}$表示第$i$个新增通道在时间步$t$处的真实EEG信号值，作为监督学习的目标标签；$\hat{\mathbf{Y}}_{i,t} \in \mathbb{R}$表示第$i$个新增通道在时间步$t$处的模型预测EEG信号值，由MASER网络生成；$C_{\text{up}} = C_h - C_l$表示超分辨率过程中新增的通道数量，即需要从低分辨率信号中预测生成的通道数目；$T \in \mathbb{N}^+$表示信号的时间步总数，对应于信号的时间长度；$\mathbf{W}_j \in \mathbb{R}^{d_j}$表示MASER网络中第$j$个可学习参数张量，包括权重矩阵和偏置向量；$P \in \mathbb{N}^+$表示网络中可学习参数总数；$\alpha \in \mathbb{R}^+$表示时域平滑约束项的权重系数，用于控制时间连续性约束的强度；$\beta \in \mathbb{R}^+$表示正则化项的权重系数，用于控制模型复杂度约束的强度。

均方误差损失项$\mathcal{L}_{\text{MSE}}$作为损失函数的主体部分，量化了模型预测信号与真实目标信号之间的重建精度。该损失项的设计基于最小二乘原理，通过最小化预测值与真实值之间的平方误差来驱动模型学习从低分辨率到高分辨率信号的准确映射关系。对于EEG信号而言，保持幅值的精确性具有特殊的重要意义，因为EEG信号的微伏级电位变化直接反映了神经元群体的活动强度和同步化程度，即使微小的幅值偏差也可能导致对神经活动模式的错误解释。该损失项通过逐点比较的方式确保重建信号在每个时间点和每个通道上都能与真实信号保持高度的数值一致性。由于其梯度计算简单且连续，有利于基于梯度的优化算法的收敛。

时域平滑约束项$\mathcal{L}_{\text{smooth}}$用于维护EEG信号在时间维度上的连续性和平滑性特征，这是生物电信号的重要生理特性之一。由于EEG信号源于神经元的生理活动，其时间演化过程受到神经系统内在动力学规律的约束，通常表现出相对平滑的变化趋势。急剧的信号跳跃或不连续现象往往对应于测量噪声、电极伪影或者非生理性的信号干扰。在超分辨率重建过程中强制保持时域平滑性不仅有助于提高信号的视觉质量，更重要的是确保重建信号符合神经生理学的基本原理。该约束项通过计算相邻时间点之间的差分平方和来量化信号的时域变化剧烈程度，差分值越大表明信号的时间变化越剧烈，平滑性越差。权重系数$\alpha$的选择需要在保持信号细节信息和确保时域平滑性之间取得平衡：过小的$\alpha$值可能导致重建信号出现不自然的高频振荡和噪声，而过大的$\alpha$值则可能过度抑制信号的正常变化，丢失重要的神经活动特征。通过在验证集上的网格搜索实验和交叉验证分析，$\alpha = 0.001$被确定为在多数EEG数据集上能够取得最优平衡的参数设置。

模型正则化项$\mathcal{L}_{\text{reg}}$采用$L_2$正则化策略，旨在通过约束模型参数的规模来防止过拟合现象的发生，并提高模型在未见数据上的泛化能力。深度神经网络由于其强大的表达能力和大量的可学习参数，在训练数据有限的情况下容易出现过拟合问题，即模型在训练集上表现优异但在测试集上性能下降。$L_2$正则化通过在损失函数中添加参数范数的平方项，鼓励模型学习较小的参数值，从而降低模型的复杂度并提高其泛化性能。正则化系数$\beta$的设置需要根据具体的网络结构和数据集特性进行调节，通常设置为较小的正值以避免过度约束模型的学习能力。在MASER的实验中，$\beta = 10^{-5}$被确定为能够有效防止过拟合同时不显著影响模型性能的最优设置。

多组件损失函数中各项权重的平衡是确保整体优化效果的关键因素。不同损失项之间可能存在相互竞争的关系，例如过度强调平滑性可能会损害重建精度，而过强的正则化可能会限制模型的表达能力。因此，权重系数的选择需要通过系统性的实验来确定最优配置。在MASER的开发过程中，采用了多阶段的超参数优化策略：首先固定$\alpha = 0$和$\beta = 0$，仅使用MSE损失训练基线模型；然后逐步引入平滑约束项，通过网格搜索确定$\alpha$的最优值；最后加入正则化项并优化$\beta$的设置。这种渐进式的权重调节策略有助于理解各损失项的独立贡献和相互作用，确保最终的损失函数配置能够充分发挥各组件的协同效应。


\section{实验设置与性能评估}
\subsection{数据集与评估指标设置}
本研究在SEED数据集\cite{zheng2015investigating}和PhysioNet MM/I数据集\cite{schalk2004bci2000, goldberger_physiobank_2000}上验证了所提出的MASER模型的性能。SEED数据集包含15名受试者在情绪相关任务期间的EEG记录，采用62个通道，采样率为1000 Hz。数据预处理过程中使用0.5至50 Hz的带通滤波器进行滤波。PhysioNet MM/I数据集包含109名受试者在运动执行和运动想象任务期间的EEG记录，采用64个通道，采样率为160 Hz。该数据集的预处理使用0.5至40 Hz的带通滤波器。
针对SEED数据集和PhysioNet数据集，本研究在三种超分辨率设置下评估MASER的性能，分别为2倍、4倍和8倍超分辨率。对于每个超分辨率因子，考虑四种不同的电极分布情况，分别为Case1、Case2、Case3和Case4。如图~\ref{fig:C3-4}和图\ref{fig:C3-5}所示，这四种情况分别对应不同的电极空间分布模式，具体电极配置在文献\cite{tang2022deep}中有详细说明。


\begin{figure}[htbp]
\centering
\includegraphics[width=\textwidth]{C3-4SEED.pdf}
\caption{SEED数据集上四种电极采样分布情况\cite{tang2022deep}}
\label{fig:C3-4}
\end{figure}

\begin{figure}[htbp]
\centering
\includegraphics[width=\textwidth]{C3-5PhysioNet.pdf}
\caption{PhysioNet MM/I数据集上四种电极采样分布情况\cite{tang2022deep}}
\label{fig:C3-5}
\end{figure}

低分辨率EEG信号通过根据Case1至Case4指定的电极配置对原始高分辨率信号进行下采样获得。实验设计采用四折交叉验证策略，对于每种情况，使用不同的随机种子进行四次独立测试。在每次测试中，数据集按7:1:2的比例划分为训练集、验证集和测试集。最终性能报告为四次运行在测试集上的平均结果及其标准差。下采样信号作为MASER的输入，然后重建高分辨率EEG信号。重建信号与原始信号进行比较，以评估所提出模型在空间分辨率增强方面的性能。这种实验设计确保了对模型性能的全面和公正评估，同时考虑了不同电极配置和超分辨率倍数对模型表现的影响。

本研究主要使用归一化均方误差（Normalized Mean Squared Error，nMSE）和皮尔逊相关系数（Pearson Correlation Coefficient，PCC）来定量评估MASER的超分辨率性能。nMSE测量预测的高分辨率信号与真实值之间的平均平方差，并通过真实值进行归一化。其定义如下：
\begin{equation}
\text{nMSE} = \frac{1}{n} \sum_{i=1}^n \left( \frac{\hat{\mathbf{Y}}_i - \mathbf{Y}_i}{\mathbf{Y}_i} \right)^2,
\end{equation}
式中，$\hat{\mathbf{Y}}_i$和$\mathbf{Y}_i$分别表示第$i$个样本的预测高分辨率信号和真实高分辨率信号，$n$表示样本总数。PCC测量预测信号与真实信号之间的线性相关性。其定义如下：
\begin{equation}
\text{PCC} = \frac{\sum_{i=1}^n (\hat{\mathbf{Y}}_i - \hat{\mathbf{Y}}_{\text{avg}})(\mathbf{Y}_i - \mathbf{Y}_{\text{avg}})}{\sqrt{\sum_{i=1}^n (\hat{\mathbf{Y}}_i - \hat{\mathbf{Y}}_{\text{avg}})^2 \sum_{i=1}^n (\mathbf{Y}_i - \mathbf{Y}_{\text{avg}})^2}},
\end{equation}
式中，$\hat{\mathbf{Y}}_{\text{avg}}$和$\mathbf{Y}_{\text{avg}}$分别表示预测信号和真实信号的均值。

nMSE通过量化预测信号与实际真实值的偏差来提供预测精度的度量。较低的nMSE值表示更好的重建质量。PCC提供模型捕获EEG数据潜在模式能力的洞察。较高的PCC值表示预测信号与真实信号之间具有更强的线性相关性，反映了模型在保持原始信号时间动态特征方面的能力。这两个指标的结合使用能够全面评估超分辨率方法在精度和相关性方面的性能。

\subsection{超分辨率性能对比与分析}
本研究提出的MASER方法是首个将状态空间模型应用于EEG超分辨率任务的方法，通过多分辨率状态空间模型实现对EEG信号的有效建模。为了全面评估MASER的性能，本研究选择了具有代表性的基线方法进行对比分析。其中，统计插值方法包括球面插值算法SI \cite{freeden1984spherical}，该方法基于电极的空间位置进行传统的数学插值。基于CNN的上采样方法包括EEGSR-GAN \cite{corley2018deep}、Deep-CNN \cite{han2018feasibility}和Deep-EEGSR \cite{tang2022deep}，这些方法利用卷积神经网络的特征提取能力和生成对抗网络的优化策略来实现超分辨率重建。基于Transformer的方法包括受掩码自编码器启发的ESTFormer和ESTFormer+CNN \cite{li2023estformer}，这些方法通过自注意力机制建模全局时空依赖关系。MASER与现有方法的主要区别在于其独特的状态空间建模框架，能够在保持线性计算复杂度的同时有效捕获长程时空依赖关系。

\begin{table*}[htbp]
\centering
\caption{脑电空间超分辨率方法的性能比较}
\resizebox{1.0\textwidth}{!}{
\begin{tabular}{cccccccc}
\toprule
\multirow{2}{*}{数据集} & \multirow{2}{*}{方法} & \multicolumn{2}{c}{缩放因子 2} & \multicolumn{2}{c}{缩放因子 4} & \multicolumn{2}{c}{缩放因子 8} \\
\cmidrule(lr){3-4} \cmidrule(lr){5-6} \cmidrule(lr){7-8}
 &  & nMSE ($\downarrow$) & PCC ($\uparrow$) & nMSE ($\downarrow$) & PCC ($\uparrow$) & nMSE ($\downarrow$) & PCC ($\uparrow$) \\
\midrule
\multirow{7}{*}{SEED} & SI \cite{freeden1984spherical} & 0.684/0.000 & 0.585/0.000 & 0.703/0.000 & 0.567/0.000 & 0.755/0.000 & 0.527/0.000 \\
 & EEGSR-GAN \cite{corley2018deep} & 0.514/0.002 & 0.693/0.001 & 0.597/0.002 & 0.629/0.001 & 0.652/0.001 & 0.583/0.001 \\
 & Deep-CNN \cite{han2018feasibility} & 0.485/0.005 & 0.713/0.004 & 0.608/0.006 & 0.622/0.006 & 0.672/0.007 & 0.569/0.002 \\
 & Deep-EEGSR \cite{tang2022deep} & 0.413/0.010 & 0.762/0.007 & 0.524/0.002 & 0.682/0.002 & 0.580/0.005 & 0.641/0.003 \\
 & ESTFormer \cite{li2023estformer} & 0.315/0.014 & 0.827/0.009 & 0.353/0.007 & 0.803/0.004 & 0.425/0.011 & 0.757/0.007 \\
 & ESTFormer+CNN \cite{li2023estformer} & \uline{0.277/0.014} & \uline{0.852/0.008} & \uline{0.311/0.005} & \uline{0.829/0.003} & \uline{0.377/0.009} & \uline{0.788/0.005} \\
 & MASER & \textbf{0.274/0.032} &\textbf{0.877/0.007} & \textbf{0.306/0.019} & \textbf{0.829/0.006} & \textbf{0.319/0.024} & \textbf{0.795/0.004} \\
\midrule
\multirow{7}{*}{\parbox{1.5cm}{PhysioNet\\MM/I}} & SI \cite{freeden1984spherical} & 0.287/0.000 & 0.855/0.000 & 0.317/0.000 & 0.839/0.000 & 0.367/0.000 & 0.811/0.000 \\
 & EEGSR-GAN \cite{corley2018deep} & 0.224/0.000 & 0.887/0.001 & 0.258/0.002 & 0.860/0.001 & 0.308/0.003 & 0.832/0.003 \\
 & Deep-CNN \cite{han2018feasibility} & 0.206/0.019 & 0.900/0.007 & 0.231/0.021 & 0.879/0.012 & 0.276/0.044 & 0.855/0.019 \\
 & Deep-EEGSR \cite{tang2022deep} & 0.165/0.001 & \uline{0.920/0.000} & 0.177/0.001 & 0.907/0.000 & 0.214/0.002 & 0.885/0.001 \\
 & ESTFormer \cite{li2023estformer} & 0.174/0.004 & 0.908/0.002 & 0.189/0.002 & 0.900/0.001 & 0.237/0.007 & 0.872/0.004 \\
 & ESTFormer+CNN \cite{li2023estformer} & \uline{0.155/0.003} & 0.918/0.002 & \uline{0.166/0.002} & \uline{0.912/0.001} & \uline{0.207/0.006} & \uline{0.890/0.004} \\
 & MASER & \textbf{0.109/0.004} & \textbf{0.960/0.006} & \textbf{0.131/0.005} & \textbf{0.928/0.004} & \textbf{0.155/0.006} & \textbf{0.896/0.005} \\
\bottomrule
\end{tabular}}
\begin{tablenotes}
\footnotesize
\item[]\centering 注：最佳结果用\textbf{粗体}标示，次佳结果用\underline{下划线}标示。nMSE和PCC指标报告为均值/标准差。
\end{tablenotes}
\label{TC3-1}
\end{table*}

表~\ref{TC3-1}中呈现的实验结果证明了MASER在各种数据集和超分辨率尺度上的优越性能。在SEED数据集上，MASER在所有缩放因子下都表现出显著的性能提升。在2倍缩放因子下，MASER相比表现最优的基线方法ESTFormer+CNN在nMSE上减少了0.8\%，在PCC上提高了2.9\%，展现出较为稳定的重建精度。在4倍缩放因子下，MASER的优势更加明显，相比ESTFormer+CNN在nMSE上减少了6.7\%，在PCC上提高了1.9\%。值得注意的是，在8倍缩放因子这一极具挑战性的高倍超分辨率任务中，MASER展现出了最为突出的性能优势，相比ESTFormer+CNN在nMSE上实现了12.7\%的显著减少，在PCC上提高了1.8\%。在PhysioNet数据集上观察到更加显著的性能提升趋势，MASER在所有缩放因子下始终获得最低的nMSE和最高的PCC值。在2倍缩放因子下，MASER相比表现最优的基线方法ESTFormer+CNN将nMSE减少了28.2\%，PCC提高了4.6\%，展现出卓越的重建质量。在4倍缩放因子下，MASER相比ESTFormer+CNN将nMSE减少了25.4\%，PCC提高了1.7\%，保持了稳定的性能优势。在8倍缩放因子的极限挑战场景下，MASER相比ESTFormer+CNN实现了最为显著的性能跃升，nMSE减少了25.1\%，PCC提高了6.7\%，充分证明了所提方法在处理复杂超分辨率任务时的卓越能力。

对实验结果分析可知，MASER的性能优势随着超分辨率倍数的增加而愈发明显，这一现象充分证明了状态空间模型在处理长序列EEG数据方面的独特优势。传统的CNN方法在高倍超分辨率任务中往往受限于感受野大小，难以建模长程时空依赖关系。而基于Transformer的方法虽然具有全局建模能力，但其二次计算复杂度在处理长序列时面临效率瓶颈。相比之下，MASER通过状态空间模型的线性复杂度设计和eSS2D模块的多方向扫描机制，能够有效捕获EEG信号的复杂时空相关性，特别是在高倍数超分辨率任务中展现出显著的性能优势。

\subsection{消融实验与参数设置}
本研究开展了全面的消融实验，旨在系统分析MASER架构中关键模块的作用机制，并确定重要参数的最优配置。通过控制变量的实验设计，深入探讨了不同编码解码组合、特征提取块中状态空间模型参数以及损失函数各组件对模型性能的具体影响，为MASER的设计合理性提供实证支撑。

（一）MASER架构中关键模块的影响

为全面评估MASER中编码器和解码器模块配置的有效性，本研究设计了七种不同的架构组合进行对比分析，如图~\ref{fig:C3-6}所示。这七种配置包括：eMamba-eMamba（eM-eM，即本研究提出的MASER采用的完整配置）、Transformer-Transformer（T-T）、Mamba-Transformer（M-T）、Transformer-Mamba（T-M）、Mamba-Mamba（M-M）、eMamba-Transformer（eM-T）和Transformer-eMamba（T-eM）。其中，eMamba是本章针对EEG信号特性专门设计的改进模块，而其他组件均为现有的成熟架构。

\begin{figure}[htbp]
	\centering
        \includegraphics[width=0.9\textwidth]{C3-6消融.pdf}
	\caption{MASER中七种不同编码解码器组合配置}
        \caption*{\centering (a)T-T组合；(b)M-T或eM-T组合；(c)T-M或T-eM组合；(d)M-M或eM-eM组合。其中T表示Transformer，M表示Mamba，eM表示eMamba。}
\label{fig:C3-6}
\end{figure}

表~\ref{TC3-2}展示了不同架构组合在SEED和PhysioNet MM/I数据集上的性能对比结果。实验结果显示，eM-eM配置实现了最佳性能，充分验证了eMamba模块设计的有效性。具体而言，MASER在SEED数据集上实现了0.290的nMSE和0.833的PCC，在PhysioNet数据集上实现了0.145的nMSE和0.925的PCC，均为所有配置中的最优结果。与次佳配置eMamba-Transformer相比，MASER在SEED数据集上将nMSE减少了3.7\%，PCC提高了0.4\%；在PhysioNet数据集上，MASER将nMSE减少了22.0\%，PCC提高了2.9\%。
通过对比分析不同配置的性能表现，可以观察到使用eMamba作为编码器的配置始终优于使用Transformer编码器的对应配置。例如，eMamba-Transformer在SEED数据集上实现了0.301的nMSE和0.830的PCC，均优于Transformer-Transformer的0.322 nMSE和0.789 PCC。这一现象表明eMamba在从低分辨率EEG信号中提取相关特征方面具有更强的能力，能够更好地捕获EEG信号的时空依赖关系和生理特性。

\begin{table}[htbp]
\centering
\caption{不同编码解码器组合的性能对比}
\begin{tabular}{cccccc}
\toprule
\multirow{2}{*}{编码器} & \multirow{2}{*}{解码器} & \multicolumn{2}{c}{SEED} & \multicolumn{2}{c}{PhysioNet MM/I} \\
\cmidrule(lr){3-4} \cmidrule(lr){5-6}
 &  & nMSE ($\downarrow$) & PCC ($\uparrow$) & nMSE ($\downarrow$) & PCC ($\uparrow$) \\
\midrule
Transformer（T） & Transformer（T） & 0.322 & 0.789 & 0.317 & 0.813 \\
Mamba（M） & Transformer（T） & 0.318 & 0.801 & 0.190 & 0.897 \\
Transformer（T） & Mamba（M） & 0.319 & 0.792 & 0.213 & 0.884 \\
Mamba（M） & Mamba（M） & 0.327 & 0.747 & 0.193 & 0.896 \\ \midrule
eMamba（eM） & Transformer（T） & 0.301 & \uline{0.830} & 0.186 & \uline{0.899} \\
Transformer（T） & eMamba（eM） & \uline{0.301} & 0.743 & \uline{0.174} & 0.896 \\
eMamba（eM） & eMamba（eM） & \textbf{0.290} & \textbf{0.833} & \textbf{0.145} & \textbf{0.925} \\
\bottomrule
\end{tabular}
\begin{tablenotes}
\footnotesize
\item[]\centering 注：最佳结果用\textbf{粗体}标示，次佳结果用\underline{下划线}标示。
\end{tablenotes}
\label{TC3-2}
\end{table}

为进一步定性分析不同架构配置的重建效果，选择四种具有代表性的配置：T-T、eM-T、T-eM和eM-eM进行深入比较。T-T代表传统的Transformer架构，eM-eM代表本章提出的完整MASER架构，而eM-T和T-eM则分别体现编码器和解码器中eMamba模块的独立作用。图\ref{fig:C3-7}展示了这些配置在时域重建方面的性能比较。结果表明，MASER生成的波形在视觉上与真实信号最为相似，展现出高水平的时域保真度。这一优势在PhysioNet MM/I数据集上尤为明显，MASER的重建结果在捕获显著峰值和谷值特征方面明显优于其他配置。图\ref{fig:C3-8}进一步展示了不同方法在空间重建方面的性能差异。SEED数据集的地形图清晰地显示，MASER的重建结果在不同空间区域和尺度上与真实信号保持更高的一致性，能够更准确地保持正确的空间模式和强度分布。特别是在红色虚线框突出显示的关键区域，MASER表现出明显的优势，证明其在捕获EEG信号空间分布特征方面的有效性。


\begin{figure*}[htbp]
	\centering
		\includegraphics[width=\textwidth]{C3-7时域图.pdf}
	\caption{不同架构配置的时域重建效果对比}
\label{fig:C3-7}
\end{figure*}


\begin{figure*}[htbp]
	\centering
	\includegraphics[width=\textwidth]{C3-8地形图.pdf}
	\caption{不同架构配置的空间重建效果对比}
    \caption*{\centering 注：第一行（GT）显示真实信号的地形图，后续行分别显示T-T、eM-T、T-eM和eM-eM（MASER）方法生成的地形图。MASER产生的地形图在视觉上与真实值最接近，特别是在红色虚线框突出显示的区域。}
\label{fig:C3-8}
\end{figure*}

（二）SSMs的关键参数设置

为确定eMamba块中SSMs的最优参数配置，本研究深入探讨了两个关键参数对模型性能的影响：深度参数$d$和状态维度$d_{\text{state}}$。深度参数$d$控制状态空间模型的层次结构复杂度，而状态维度$d_{\text{state}}$决定了模型内部状态表示的容量。通过系统性的参数组合实验，测试了六种不同的参数设置：$(1, 4)$、$(1, 8)$、$(1, 16)$、$(2, 4)$、$(2, 8)$和$(2, 16)$，以探索参数配置对模型性能的具体影响规律。


\begin{table}[htbp]
\centering
\caption{不同状态空间模型参数配置的性能对比}
\begin{tabular}{cccccc}
\toprule
\multirow{2}{*}{$d$} & \multirow{2}{*}{$d_{state}$} & \multicolumn{2}{c}{SEED} & \multicolumn{2}{c}{PhysioNet MM/I} \\
\cmidrule(lr){3-4} \cmidrule(lr){5-6}
 &  & nMSE ($\downarrow$) & PCC ($\uparrow$) & nMSE ($\downarrow$) & PCC ($\uparrow$) \\
\midrule
1 & 4 & 0.348 & 0.816 & 0.164 & 0.915 \\
1 & 8 & 0.346 & 0.814 & 0.164 & 0.912 \\
1 & 16 & \uline{0.304} & 0.827 & 0.143 & \uline{0.927} \\
2 & 4 & 0.313 & 0.834 & \uline{0.149} & \textbf{0.928} \\
2 & 8 & \textbf{0.297} & \textbf{0.844} & \textbf{0.135} & 0.917 \\
2 & 16 & 0.309 & \uline{0.838} & 0.152 & 0.922 \\
\bottomrule
\end{tabular}
\begin{tablenotes}
\footnotesize
\item[]\centering 注：最佳结果用\textbf{粗体}标示，次佳结果用\underline{下划线}标示。
\end{tablenotes}
\label{TC3-3}
\end{table}

表~\ref{TC3-3}的实验结果表明，增加深度参数和状态维度通常能够显著改善模型性能，但当参数规模达到一定程度时，性能提升趋于饱和，甚至可能出现轻微下降的现象。这种现象反映了模型复杂度与性能之间的非线性关系，以及计算资源和参数效率之间的权衡。
具体分析各参数配置的性能表现，对于SEED数据集，$d=2$和$d_{\text{state}}=8$的配置取得了最佳的综合性能，实现了0.297的nMSE和0.844的PCC。对于PhysioNet MM/I数据集，不同参数配置在不同指标上展现出差异化的优势：$d=2$和$d_{\text{state}}=4$的配置在PCC指标上达到最优值0.928，而$d=2$和$d_{\text{state}}=8$的配置在nMSE指标上实现最低值0.135。将最优配置（$d=2$和$d_{\text{state}}=8$）与基础配置（$d=1$和$d_{\text{state}}=4$）进行对比，在SEED数据集上nMSE降低了14.7\%，PCC提高了3.4\%；在PhysioNet MM/I数据集上nMSE降低了17.7\%，显示出显著的性能提升效果。适度增加模型复杂度能够有效提升特征表示能力和重建精度，但过度复杂的参数配置可能导致过拟合和计算效率下降。$d=2$和$d_{\text{state}}=8$的配置在计算效率和模型性能之间实现了最佳平衡，成为EEG超分辨率任务的推荐参数设置。


（三）损失函数组件的消融分析

为深入理解多组件损失函数中各项的具体贡献，本研究开展了系统性的消融实验。首先通过参数搜索确定各损失组件的基线权重值，采用粗网格搜索策略，将平滑约束系数$\alpha$设置为0.001，正则化系数$\beta$设置为0.00001。随后通过控制变量法，分别测试不同损失组件组合对模型性能的影响。


\begin{table}[htb]
\centering
\caption{损失函数组件的消融研究}
\label{tab:TC3-4}
\begin{tabular}{cccccccc}
\toprule
\multirow{2}{*}{MSE} & \multirow{2}{*}{Smooth} & \multirow{2}{*}{Reg} & \multicolumn{2}{c}{SEED} & \multicolumn{2}{c}{PhysioNet MM/I} \\
\cmidrule(lr){4-5} \cmidrule(lr){6-7}
 &  &  & nMSE ($\downarrow$) & PCC ($\uparrow$) & nMSE ($\downarrow$) & PCC ($\uparrow$) \\
\midrule
$\checkmark$ &   &   & 0.325 & 0.807 & 0.181 & 0.893 \\
$\checkmark$ & $\checkmark$ &   & \uline{0.300} & \uline{0.825} & \uline{0.152} & \uline{0.920} \\
$\checkmark$ &   & $\checkmark$ & 0.314 & 0.810 & 0.163 & 0.908 \\
$\checkmark$ & $\checkmark$ & $\checkmark$ & \textbf{0.293} & \textbf{0.836} & \textbf{0.146} & \textbf{0.922} \\
\bottomrule
\end{tabular}
\begin{tablenotes}
\footnotesize
\item[]\centering 注：最佳结果用\textbf{粗体}标示，次佳结果用\underline{下划线}标示。
\end{tablenotes}
\end{table}

表~\ref{tab:TC3-4}的消融实验结果表明，在基础MSE损失的基础上引入额外约束组件能够在两个数据集上产生显著的性能改善。平滑约束的引入带来了最为显著的性能提升效果，在SEED数据集上使nMSE减少了7.7\%，PCC提高了2.2\%；在PhysioNet MM/I数据集上使nMSE减少了16.0\%，PCC提高了3.0\%。这一显著改善充分证明了平滑约束在捕获EEG信号时域连续性特征方面的有效性，有助于确保重建信号符合神经生理学的基本规律。正则化项虽然贡献程度相对较小，但仍对模型性能产生积极影响。通过对比仅使用MSE损失和同时使用MSE与正则化损失的结果，可以观察到正则化项有效缓解了过拟合现象，提升了模型的泛化能力。当同时使用MSE损失、平滑约束和正则化三个组件时，模型在两个数据集上均实现了最优性能。这些显著的性能提升验证了多组件损失函数设计的合理性，证明了不同约束项之间的协同作用能够从多个维度优化重建信号的质量，在重建精度、时域连续性和模型泛化能力之间实现了良好的平衡。消融实验的结果进一步强调了特定领域约束和正则化技术在EEG超分辨率任务中的重要性，为深度学习方法在神经信号处理领域的应用提供了有价值的设计指导。

\subsection{模型计算效率分析}
从复杂度角度分析，eMamba块展现出显著的计算优势。eMamba块的计算复杂度为$O(N \times D)$，其中$N$表示输入序列的长度，$D$表示状态空间模型的维度。这种线性复杂度的特性源于状态空间模型的递归计算结构，每个时间步的状态更新仅依赖于前一时间步的状态和当前输入，避免了全局依赖关系的计算。相比之下，Transformer架构中的自注意力机制需要计算序列中每对元素之间的相关性，导致其计算复杂度为$O(N^2 \times D)$。这种二次复杂度的增长模式在处理长序列数据时会带来显著的计算负担。具体而言，当序列长度从1000增长到10000时，Transformer的计算量增长100倍，而eMamba的计算量仅增长10倍，展现出明显的扩展性优势。
计算效率的提升对EEG超分辨率任务具有重要的实际意义。首先，高效的计算使得模型能够处理更长时间窗口的EEG数据，有助于捕获更完整的神经活动模式和长程时间依赖关系。其次，较低的计算资源需求降低了模型部署的硬件门槛，使得MASER能够在资源受限的环境中得到应用。最后，高效的推理速度为实时EEG分析和在线BCI应用提供了技术基础。


\section{运动想象识别案例研究}

\subsection{案例与方法设置}
为验证MASER在实际BCI应用中的有效性，本研究设计了基于PhysioNet MM/I数据集的运动想象识别案例研究。该案例的核心目标是探索通过算法手段提升空间分辨率是否能够有效改善运动想象识别性能，从而为消费级EEG设备在实际应用中的性能提升提供技术支撑。研究聚焦于左手和右手运动想象任务的二分类问题，采用广泛应用的EEGNet\cite{lawhernEEGnet2018}作为分类模型，通过比较不同分辨率配置下的分类性能，评估MASER算法在空间分辨率补充方面的有效性。

本研究设计了多组对照实验以全面评估不同空间分辨率配置下的分类性能。实验组包括：使用16通道、32通道和64通道真实EEG信号进行运动想象分类的基准组，以及使用MASER将16通道信号超分辨率增强至32通道和64通道的实验组。通过系统比较这些配置下获得的分类准确率和kappa系数，量化分析MASER在空间分辨率增强方面的贡献，并验证高分辨率EEG信号在运动想象分类任务中的潜在优势。

\subsection{分类性能对比分析}

表~\ref{tab:TC3-5}系统展示了不同分辨率配置下EEG信号在运动想象二分类任务中的性能对比结果。通过对比分析不同配置的性能表现，可以观察到明显的分辨率效应。16通道基线信号的分类准确率为83.2\%，kappa系数为0.625，代表了低分辨率配置的性能基准。随着空间分辨率的提升，32通道真实信号的准确率提升至84.8\%，kappa系数达到0.640，显示出适度的性能改善。更为显著的性能提升出现在64通道配置中，无论是真实信号还是超分辨率信号，都实现了接近88\%的高准确率，相较于16通道基线提升了约5\%。值得关注的是，MASER生成的超分辨率信号在多个方面展现出优异的性能特征。32通道超分辨率信号实现了84.2\%的准确率和0.631的kappa系数，与对应的真实信号性能极为接近，验证了超分辨率技术在中等分辨率增强方面的有效性。64通道超分辨率信号的表现更为突出，其kappa系数甚至略优于真实64通道信号，这一现象表明超分辨率过程不仅能够有效恢复空间信息，还可能通过隐含的去噪和特征增强机制改善信号质量。

\begin{table}[htb]
\centering
\caption{不同分辨率配置下EEG信号运动想象分类性能对比}
\begin{tabular}{ccc}
\toprule
通道配置 & 准确率 (\%, $\uparrow$) & Kappa ($\uparrow$) \\
\midrule
16通道 (Base, GT) & 83.2±3.9 & 0.625±0.064 \\  \midrule
32通道 (GT) & 84.8±2.8 & 0.640±0.052 \\
32通道 (SR) & 84.2±2.5 & 0.631±0.058 \\ \midrule
64通道 (GT) & \textbf{88.1±2.2} & \uline{0.733±0.058} \\
64通道 (SR) & \uline{88.0±2.2} & \textbf{0.735±0.047} \\
\bottomrule
\end{tabular}
\begin{tablenotes}
\footnotesize
\item[]\centering 注：Base表示基线分辨率配置，GT表示真实高分辨率信号，SR表示使用MASER生成的超分辨率信号。最佳结果用\textbf{粗体}标示，次佳结果用\underline{下划线}标示。
\end{tablenotes}
\label{tab:TC3-5}
\end{table}

实验结果充分证实了MASER在增强低分辨率EEG系统空间分辨率方面的技术优势，生成的高分辨率信号能够达到与真实高分辨率信号相当的分类性能水平。该发现表明通过先进的信号处理算法可以显著突破硬件配置的限制，为成本受限的便携式BCI系统提供性能提升的技术路径。



\subsection{特征空间可视化分析}
为深入理解超分辨率技术对运动想象分类任务的影响机制，本研究采用主成分分析（Principal Component Analysis, PCA）方法对不同分辨率配置下提取的EEG特征进行降维可视化分析。图\ref{fig:C3-9}展示了EEGNet分类网络最终特征层的二维主成分投影结果，通过特征空间的几何分布特性定性评估不同配置下左手和右手运动想象任务的可分离性。


\begin{figure}[htb]
	\centering
		\includegraphics[width=\textwidth]{C3-9PCA-CN.pdf}
    \caption{EEGNet特征层的主成分分析可视化结果}
    \caption*{\centering （a）16通道真实信号；（b）16通道真实信号的特征PCA效果；（c）32通道真实信号的特征PCA效果；（d）32通道超分辨率信号的特征PCA效果；（e）64通道真实信号的特征PCA效果；（f）64通道超分辨率信号的特征PCA效果。注：不同颜色点分别代表左手和右手运动想象任务的特征分布。}
\label{fig:C3-9}
\end{figure}

主成分分析结果直观地展现了空间分辨率对特征可分离性的显著影响。从可视化结果可以清晰观察到，随着通道数量的增加，不同运动想象类别在特征空间中的聚类紧密度和类间分离度均呈现出明显的改善趋势。16通道配置下的特征分布相对分散，类别边界较为模糊，反映了低分辨率配置在特征表示能力方面的局限性。32通道配置显示出适度的改善，特征聚类结构更为清晰，类间分离度有所提升。64通道配置下的可视化结果最为理想，无论是真实信号还是超分辨率信号，都展现出优异的特征聚类特性和类间分离效果。左手和右手运动想象任务在特征空间中形成了相对紧密且分离良好的聚类结构，为分类算法提供了清晰的决策边界。MASER生成的64通道超分辨率信号的特征分布与真实64通道信号极为相似，两者在主成分空间中展现出几乎一致的聚类模式和分离特性。


综合主成分分析的定性结果和分类性能的定量指标，可以得出一致的结论：MASER技术能够有效增强低分辨率EEG系统的空间分辨率，生成的超分辨率信号在特征表示和分类性能方面均达到了与真实高分辨率信号相当的水平。这一研究发现验证了超分辨率技术在BCI应用中的技术可行性，为消费级EEG设备通过算法手段实现性能突破提供了有力的实践指导。


\section{本章小结}
本章针对EEG信号空间分辨率不足这一制约EEG技术广泛应用的关键瓶颈，提出了基于状态空间模型的超分辨率算法MASER。该方法将从低密度到高密度EEG信号的重建过程建模为条件生成问题，通过融合卷积操作的局部结构感知能力与状态空间模型的长程依赖建模优势，设计了eMamba模块作为核心组件。MASER采用编码器--解码器架构，结合eSS2D模块的多方向选择性扫描机制，以线性计算复杂度有效捕获EEG信号的复杂时空依赖关系，并通过多组件复合损失函数确保重建信号的生理合理性和时域连续性。
实验结果表明，MASER在SEED和PhysioNet MM/I数据集上的多种超分辨率设置下均显著优于传统插值、CNN和Transformer等基线方法，特别是在高倍数超分辨率任务中展现出更加突出的性能优势。消融研究验证了eMamba模块设计的有效性以及各损失组件的协同作用。运动想象识别案例进一步证实，MASER生成的超分辨率信号能够达到与真实高分辨率信号相当的分类性能，为消费级BCI系统通过算法途径突破硬件限制、实现“软件补偿硬件”的空间分辨率提升提供了可行的技术路径。MASER在空间超分辨率重建方面取得了显著进展，但电极偏移、阻抗变化等实际部署挑战以及跨会话适应性问题将是未来工作的重要方向。

